%% sec_conclusion.tex
%% ─────────────────────────────────────────────────────────────────────────────
%% Section 7 — Conclusion and Future Work
%% ─────────────────────────────────────────────────────────────────────────────

\section{Conclusion}\label{sec:conclusion}

We have presented a deterministic wavefunction-token cascade architecture for
leading-order Monte Carlo matrix element evaluation on AMD Versal \aie arrays,
targeting the \ggttg process with 16 Feynman diagrams.  The central finding is
that the 16\KB program-memory limit of a single \aie tile---a hard, non-scalable
hardware constraint---fundamentally determines the pipeline structure.  No
single tile can hold the complete \ggttg computation (${\sim}38\KB$ required);
the five-kernel cascade pipeline is not an optimisation choice but a
\emph{necessary architectural response} to this constraint.

The key contributions and take-away messages are:

\begin{itemize}

  \item \textbf{Program-memory-driven partitioning is a systematic
    methodology.}  Profiling HELAS function PM footprints, identifying
    exclusion constraints (VVV Exclusion Principle), and applying lazy
    evaluation provide a reproducible recipe for mapping any
    \mgfive-generated ME computation onto an \aie array.

  \item \textbf{The cascade token protocol enables stateful multi-tile
    pipelines.}  Passing partial wavefunctions and JAMP accumulators over the
    deterministic 384-bit cascade interface eliminates the need for any
    shared memory or flow-control infrastructure, while the cascade contract
    (identical loop skeletons, unconditional writes) provides a simple and
    verifiable deadlock-freedom guarantee.

  \item \textbf{The design achieves 100\% tile utilisation.}  The five-tile
    pipeline granularity fits the $50\times 8 = 400$-tile array exactly (80
    pipelines), leaving no tiles idle and no routing resources wasted.

  \item \textbf{The system is compute-bound, not I/O-bound.}  PLIO utilisation
    is below 0.5\% in all scenarios, confirming that further speedup should
    target kernel compute latency (e.g., helicity filtering for a $2\times$
    improvement) rather than I/O optimisation.

\end{itemize}

\paragraph{Limitations.}
Only the LO \ggttg process is demonstrated.  Precision is float32
(${\sim}10^{-7}$ relative error vs.\ double-precision), which is acceptable for
LO calculations but would require attention at NLO.  Throughput numbers at
full 80-pipeline scale are currently projected from single-pipeline measurements
and require hardware validation.

\paragraph{Future work.}
\begin{itemize}
  \item \textbf{Helicity filtering.}  Pre-computing a good-helicity bitmask
    at runtime can reduce the 32-helicity inner loop to ${\sim}16$, providing
    a $2\times$ speedup within the same loop-skeleton framework.
  \item \textbf{Higher-multiplicity processes.}  The methodology generalises:
    $\gluglu\to\ttbar\gluglu\gluglu$ (178 diagrams) would require a deeper
    pipeline spanning more tiles; the token format and cascade contract extend
    without modification to the architecture.
  \item \textbf{NLO integration.}  Real-emission and loop corrections involve
    additional diagram topologies but the same HELAS function set.  The PM
    partitioning methodology applies directly; the primary challenge is the
    increased token complexity for loop-integral wavefunction products.
  \item \textbf{Direct hardware benchmark} of the full 80-pipeline array to
    confirm the linear-scaling throughput projection and measure the
    end-to-end energy per PSP.
\end{itemize}

\backmatter

\bmhead{Acknowledgements}
This work was supported by \TODO{funding source and grant number}.
The authors acknowledge the use of the VCK190 evaluation kit provided by AMD
and the computing resources of \TODO{institution}.

\section*{Declarations}

\noindent\textbf{Funding.}
\TODO{List all funding sources with full organisation names and grant numbers.
Example: This work was supported by the European Research Council (ERC) under
the European Union's Horizon~2020 programme (grant agreement No.\ xxxxxx).}

\noindent\textbf{Conflict of interest.}
The authors declare that they have no competing interests, financial or
non-financial, that are directly or indirectly related to the work submitted
for publication.

\noindent\textbf{Ethics approval and consent to participate.}
Not applicable.

\noindent\textbf{Consent for publication.}
Not applicable.

\noindent\textbf{Data availability.}
The HELAS kernel source code and benchmark datasets supporting the results
of this study are available from the corresponding author upon reasonable
request.  The AMD Versal VCK190 platform and Vitis/AIE toolchain used are
commercially available.

\noindent\textbf{Code availability.}
The AIE kernel source code and host application are available from the
corresponding author upon reasonable request.

\noindent\textbf{Author contribution.}
\TODO{Describe each author's contribution following CRediT taxonomy:
Conceptualization, Methodology, Software, Formal analysis, Investigation,
Data curation, Writing -- original draft, Writing -- review \& editing,
Visualization, Supervision, Funding acquisition.}
